\chapter{Závěr}
\label{chap:conclusion}

Cílem této práce bylo vytvořit nástroj pro sledování aktivity rozsáhlých GitHub repozitářů, konkrétně repozitáře překladače jazyka Rust (\texttt{rust-lang/rust}).
Motivací byl fakt, že běžně dostupné nástroje nedokážou pokrýt potřeby masivních open-source projektů, jejichž vývojový cyklus je řízen složitou organizační strukturou a~stavovými automaty nad štítky.

V~teoretické části byl nastudován datový model GitHub repozitářů a~dostupná \API, od \REST přes GraphQL až po Webhooky.
Tato rešerše poskytla základ pro návrh systému.

Na základě získaných znalostí byla navržena struktura relační databáze pro agregaci metadat repozitáře z~různých zdrojů: GitHub \API, lokální Git historie a~konfiguračních souborů týmů.
Výsledný systém byl implementován v~jazyce Rust, který umožňuje jednoduché znovuvyužití jako knihovny v~rámci jiných nástrojů.

Nad navrženou databází byla implementována sada analytických dotazů pro získání užitečných informací o~repozitáři.
Dotazy umožňují zjistit například to, jak dlouho čekají \PR na code review, kteří přispěvatelé pracují v~konkrétních oblastech kódu nebo jaký je vývoj počtu otevřených \PR v~čase.
Výkonnostní evaluace prokázala, že s~vhodně navrženými databázovými indexy lze tyto dotazy provádět efektivně i~nad stovkami tisíc záznamů.

Součástí implementace je rovněž interaktivní webové rozhraní, které výsledky analytických dotazů vizualizuje formou přehledných grafů a~tabulek.

Nástroj byl otestován na datech repozitáře překladače Rustu, což je netriviální projekt s~více než 200\,000 issues a~\PR.
Praktické vyhodnocení ukázalo, že nástroj dokáže identifikovat reálné bottlenecky ve schvalovacím procesu, například odlišit čas čekání na autora od čekání na recenzenta, což není pomocí standardního rozhraní GitHub možné.
Všechny body zadání tak byly splněny.

\section*{Budoucí rozvoj}
\label{sec:conclusion:rozvoj}

Nabízí se několik směrů pro budoucí rozšíření nástroje.
V~první řadě by bylo možné doplnit další typy analytických dotazů, například analýzu průměrné doby od otevření \PR po jeho sloučení v~závislosti na týmu nebo komponentě, identifikaci souborů nejčastěji stojících za konflikty při slučování nebo sledování dlouhodobých trendů v~aktivitě jednotlivých kontributorů.
Druhým přirozeným rozšířením je obohacení sledovaných dat: v~současnosti databáze neobsahuje například komentáře k~\PR, výsledky \CI nebo závislosti mezi issues.
Zahrnutí těchto informací by umožnilo hlubší analýzy a~otevřelo cestu k~prediktivním modelům, například k~odhadování rizika zanesení chyby při sloučení konkrétního \PR.

\endinput